% 本文件由 LaTeX 排版 Agent 根据有效 translation.json 自动生成。
% author=Ephraim the Syrian; work_id=3707; title=劝诫与悔改讲道

\chapter{劝诫与悔改讲道}

% structure: p0001 source=h1 target=omitted reason=page-title-already-rendered-by-parent

1. 道理不是出于强迫；生命之言出于自由意志。凡愿意听道理的，让他洁净自己意愿的田地，好让好种子不落在虚妄探询的荆棘中。你若留心生命之言，当断绝邪恶之事；听言语对忙于罪恶的人毫无益处。你若愿意成为好人，不要喜爱放荡的习俗。首先，信赖天主，然后听从祂的法律。\par

2. 你不能听祂的话，当你还不认识你自己；你若遵守祂的审判，而你的理智远离祂，谁将给你报酬？谁将为你保存赏报？你曾因祂的名受洗；宣认祂的名！在\OriginalTerm{位格}{Persons}与命名上，圣父、圣子、圣神，三个名与三个位格，这三者将成为你的屏障，抵挡分裂与争辩。不要怀疑真理，免得你因真理而丧亡。你从水中受洗；你在祂的命名中穿上了基督；主的宝座在你身上，祂的印记在你额上。当心不要成为别人的，因为你别无他主。唯有祂在仁慈中造了我们；唯有祂在十字架上救赎了我们。祂引导我们的生命；祂掌管我们的软弱；祂实现我们的复活。祂按我们的行为赏报我们。凡宣认祂、听从并遵守祂诫命的人是有福的！你，人啊，是至高天主的儿子。当心不要因你的行为激怒那良善又恩慈的父。\par

3. 你若向邻人发怒，就是向天主发怒；你若心中怀怒，你的狂妄就指向你的主。你若因嫉妒而责斥，你的一切责斥都是邪恶的。但若爱德居住在你内，你在世上就没有仇敌。你若真是和平之子，就不会激起任何人的怒气。你若公正正直，就不会亏待你的同伴。你若喜爱发怒，就去对恶人发怒，那将适合你；你若要争战，看！\OriginalTerm{撒殚}{Satan}是你的对手；你若要辱骂，就向魔鬼展示你的诅咒。你若侮辱君王的肖像，就必受杀人罪的刑罚；你若辱骂一个人，就是辱骂天主的肖像。尊敬你的邻人，看！你尊敬了天主。但你若想羞辱祂，在怒气中攻击你的邻人！\par

4. \Emph{这是第一条诫命：——\BibleQuote{你要全心、全灵、全力爱主你的天主}}\BibleRef{玛 22:37}，按你所能。你爱天主的标记，就是爱你的同伴；你若恨你的同伴，你的恨便是向着天主。因为你在天主前祈祷时发怒，便是亵渎。因为你的心也谴责你，你空费言辞；你的良心公正地判断，你的祈祷毫无益处。基督悬挂在树的高处时，曾为杀害祂的人求情；而你（是）尘土、泥巴之子，怒气随意充满你。你向弟兄怀怒，还敢祈祷吗？即使站在你这边的人，虽不是你的罪之邻，不义的污点也触及他，他的祈求不被垂听。放下怒气再祈祷；除非你想进一步激怒，约束怒气，如此你才能恳求。若他（对方）不打算在狂怒中与你相遇，就从此身内驱逐怒气，因为它被私欲所束缚。\par

5. 你有灵性的本性；灵魂是造物主的肖像；尊敬天主的肖像，通过与人和谐相处。记住死亡，不要发怒，好使你的平安不是出于勉强。只要你的生命尚存，洁净你的灵魂脱离怒气；因为若它随你进入\OriginalTerm{阴府}{Sheol}，你的路将直通\OriginalTerm{地狱}{Gehenna}。不要在心里存怒；不要在灵魂中怀忿；你无权支配你的灵魂，除非去做那善事。你是用天主的血买来的；你是因基督的苦难被救赎的；祂为你忍受死亡，好使你对罪恶死。祂的面容忍受唾污，好使你不畏惧嘲笑。祂喝醋和苦胆，好使你与怒气隔绝。祂的身上受鞭打，好使你不怕受苦。你若是祂真实的仆人，要敬畏你圣洁的主；你若是祂真实的门徒，要行走在你师傅的足迹中。忍受来自弟兄的嘲笑，好使你成为基督的同伴。不要向人显示怒气，好使你不与你的救赎主分离。\par

6. 你是人，地上的尘土，泥土，与土块同族；你是野兽族类的儿子。你若不知道你的尊荣，就要以行为而不以言语，使你的灵魂与禽兽分开。你若喜爱嘲笑，就完全像\OriginalTerm{撒殚}{Satan}；你若嘲笑你的同伴，你就是魔鬼的口舌；你若对缺陷和瑕疵，在（伤害性的）名字中取乐，撒殚虽不在受造物中，但你已经强夺了他的位置。远离此，人啊，因为它完全有害；你若愿善生，不要与嘲笑者同坐，免得你成为他罪恶与惩罚的伙伴。憎恶那完全（使人哭泣的）嘲笑，与那（使人洁净的）欢乐。你若偶然听见嘲笑者，当你无意时，用光明十字架划十字，然后像羚羊一样速速离开。撒殚居住的地方，基督决不居住；一个嘲笑邻人的人，对撒殚来说是宽敞的住所；嘲笑者的心是仇敌的宫殿。撒殚不希望添加任何其他邪恶；嘲笑足以代替一切。罪人不能以他的罪充满他的肚子或钱包。他因自己的笑声而被掠夺，却不知道也不觉察。他的伤无药可治；他的病无药可愈；他的痛无可补救；他的疮无药可医。我不愿对这样的人伸出舌头责备他；他自己的羞耻足够他；他自己的狂妄足够他。那不曾听他话的人有福了；那不认识他的人有福了。教会啊，愿你远离，不让他进入你，那撒殚的邪恶酵母！\par

7. 生命的道路是狭窄的，折磨的道路是广阔的；祈祷能将人带到天国的居所。这是完美的工作：纯洁无罪的祈祷。人的义德算不得什么。人的工作是什么？他的劳苦全是虚妄。主啊，出于祢，出于祢的恩宠，我们才能在气质上变为善良。出于祢的是义德，使我们从人成为义人。出于祢的是仁慈和恩惠，使我们从尘土成为祢的肖像。赐力量给我们的意志，使我们不沉溺于罪恶！将记忆注入我们心中，使我们每时每刻都认识祢的尊荣！将真理种植在我们思想中，使我们不因疑惑而丧亡！以祢的法律占据我们的悟性，使它不因虚妄的思念而游荡！规整我们肢体的活动，使它们不给我们带来伤害！亲近天主，撒殚就会逃离你。从你心中驱除情欲，看！你已击溃了敌人。憎恨罪恶和邪恶，撒殚就会立刻逃遁。你服侍什么罪，就是在崇拜隐秘的偶像。你喜爱什么过犯，就是在你灵魂中服侍魔鬼。每当你与弟兄争竞，撒殚就安然居住。每当你嫉妒同伴，你就使魔鬼得到安息。每当你谈论不在场之人的短处，你的舌头就为魔鬼的音乐做了竖琴。每当仇恨在你灵魂中，欺骗者的平安就大得无比。每当你喜爱咒语，你的劳苦就全是左手的行为。如果你喜爱不雅之言，你就在为魔鬼预备筵席。因为这就是偶像崇拜，即情欲（肉身）的运作。\par

8. 如果你骄傲地送礼，这并非出于天主。如果你因自己的知识而高傲，你就拒绝了天主的恩宠。如果你贫穷而骄傲，看！你的结局就在你的折磨中。如果你傲慢而匮乏，看！你的匮乏趋向你的毁灭。如果你生病而呼号，看！你的苦难充满伤害。如果你缺乏食物，而你的心却渴望财富；你的窘迫与穷人同在，但你的折磨与富人同在。如果你淫荡地观看，并贪恋邻人的妻子，看！你的份将与奸淫者同在，你的地狱将与淫乱者同在。\Emph{\Quote{让你的泉源只为你自己，从你的井中饮水。让你的泉水只属于你自己，不让别人与你同饮。}} 你要求自己的身体纯洁，如同你要求你的伴侣一样。你不愿她与你年轻时的妻子与别的男人行淫；你也不可与别人的妻子、有不同丈夫的女人行淫。愿她的玷污在你眼中可憎；完全远离它。贞洁适宜于妻子；纯洁是她的装饰；法律适宜于丈夫；公义是他头上的冠冕。不要贪恋邻人的床榻，免得别人贪恋你的床榻。在你的婚姻中保持纯洁，使你的婚姻成为圣洁。那败坏邻人妻子的人，良心责备他。凡行淫的人，因恐惧而自欺。黑暗对他来说比光明更宝贵，因为他的生活方式不纯洁。暗中犯奸淫的人时刻恐惧。奸淫者也是贼，在黑暗中破门而入。那地方本身责备他，他在那里行邪恶和恶事。他进入卧室犯罪；在黑暗中遂其意愿。时候将到，它要被揭露，他的隐秘行为要显明。你用什么眼睛向天主祈祷？你举起什么手来求饶恕？为你自己羞愧和惊惶，因为你缺乏理智。如果你的邻人看见你，你就羞愧和惊惶，更何况在洞察万有的天主面前，你岂不更当羞愧？你像你的同伴母猪，完全在泥沼中打滚。即便在看的时候，你也能犯罪，如果你的心神不警醒；在听的时候，你也能过犯，如果你不守护你的听觉。淫乱者的心因充满不洁的言语而放荡。隐藏在心念中的情欲，视觉和听觉会唤醒它。\par

9. 那渴望行淫的人穿上羞耻的衣服，好使从衣着的贪欲中，放荡得以进入并居住在他心里。不要让你的衣着成为那公然放荡者的罗网。不要诡诈地说一句话，也不要挖掘你邻人的井。不要盯着妓女看；不要被她的美貌所诱惑。她就像疯狗，甚至比它更胆大。羞耻已从她脸上消失，她不知羞耻为何物。要以唾沫迎接她，以辱骂对付她，用棍子追她如同追狗，因为她就像狗一样，可与之相比。要拒绝她话语的甜蜜，免得落入她的网中。她掏空钱袋和行囊，且获利无数。要逃离她，因为她是毒蛇之女，免得她撕裂你的全身。\par

10. 你不可毁谤任何人，免得他们称你为撒殚。如果你恨这名字，就不要接近这行为；但如果你爱这行为，就不要因这名字而发怒。首先，从野兽和飞鸟那里领受责备：它们如何各自与其同类相聚；同样，你也要与你的伴侣和睦。不要因人的羞辱而欢喜，免得你自己也成为撒殚。如果恶事临到恨你的人，你不可欢喜，免得你犯罪。如果你的仇敌跌倒，你要痛苦哀伤。你要谨守你的心，免得它暗中犯罪；因为思想和行为都要被揭露。用你的双手劳作，让你的心在祈祷中默想。不爱无益的闲谈，因为对灵魂和肉体都有益的谈话会减轻你劳苦的负担。\par

11. 穷人是否在你门口哭泣？起身欣然为他开门：当他疲倦时，使他振作；扶持他的心，因它忧伤。你亲身经历过贫困的苦楚：不可只接待他人进屋，却将乞丐赶出。你也当有一部法律，一部合宜的家规。建立明智的秩序，免得卑贱者嘲笑你。在你一切行事上谨慎，免得成为愚人的笑柄。要正直而审慎，既纯朴又明智。让你的身体安静而愉快，你的问候合宜而简朴；你的言语无可指责，你的谈论简短而甘美；你的话语少而健全，充满甘饴和悟性。不可多言，即使是有智慧的话；因为一切过多的事，虽有智慧，也令人厌倦。——对待你的家人，当如父亲。在弟兄中，当自视最小，在同伴中为卑下，在众人眼中为无足轻重。与朋友保守秘密；对爱你的人要真诚。务要避免争吵；不可泄露朋友的秘密，免得听见的人都恨你，视你为挑拨是非者。不可与恨你的人争吵，无论是当面还是心里。除了撒旦本身，你不可有任何仇敌。给你所娶的妻子提出劝告；留意她的行为；因你更强壮，理当扶持她的软弱。因为女子是软弱的，极易跌倒。当如鹰，在（怒气）被点燃时，但当怒气离开你时，要喜乐且坚定，在品质的调和之中。在长者前保持沉默；对老人给予应有的尊敬。当殷勤尊敬司铎，如同家中良善的管家。按他们的职分给予应有的尊敬，不要查究他们的行为。司铎在其职分上是天使，但在其行为上是一个人。因仁慈他成为中保，在天主与人类之间。\newline{}\par

12. 不可查究人的过错；不可揭露你同伴的罪恶；不可用口舌谈论邻人的短缺。你在受造界并非审判者，你在地上并无权柄。你若爱正义，就应责备你的灵魂和你自己。作你自己罪过的审判者，自己过犯的惩罚者。不可恶意查究人的恶行。因为若你这样做，伤害必不缺少。不可轻信耳闻之言，因为欺骗者众多。不要相信虚妄的报告，因为虚假的传闻并不少。\newline{}\par

13. 不可注重符咒和占卜，因为那是与撒旦的交往。不可喜爱空谈，即使是为了正义。不可为自己开启谈论，即使是为了合宜之事。要逃避并躲避争吵，如同躲避凶暴的强盗。不可为人作保，免得犯罪。按你所拥有的帮助比你更穷的人。不可嘲笑愚人，祈祷你也不要像他。不可责备犯罪的人，免得你也被羞辱。对悔改自己罪恶的人，要作帮助者和劝告者，鼓励能够站起来的人。让他持守对天主的望德，他的罪必如碎秸焚尽。要探望病人，不可厌倦，好得众人的喜爱。要熟悉哀伤之家，却要作宴乐之家的陌生人。不可常饮酒，免得你的过犯增多。要为你的嘴唇筑起墙垣，为你的口设立守卫。要忍受与邻人的苦难，也分担他的忧患。在忧患中，良朋益友被爱他的人认识。当以爱德追随亡者，带着哀伤和奉献，并祈祷他在所往的隐密处得安息。\newline{}\par

14. 当你站立祈祷时，要在你的灵魂中呼喊：\Quote{可怜我吧，我是罪人且软弱；天主啊，求祢仁慈对待我的软弱，赐给我力量，使我能够献上中悦祢旨意的祈祷。\InnerQuote{求祢不惩罚我的仇敌，不对恨我的人施行报复；但求祢在恩宠中赐予他们，使他们成为祢旨意的履行者。}}\newline{}\par

15. 不要抵抗恶，因为凡抵抗恶的，便是从恶者而来。不可向任何人扣留任何东西，免得他若灭亡，你受责备。不可因财产和所有物而改变对人的尊重。要视万物如无，唯独天主存在。若你向邻人有所求，他却不照你所愿给予，你不可在怒气中说一句充满苦毒的话。不可违逆〔合宜的〕节期，因为世事多变。要使忧远离你的肉体，悲伤远离你的思绪；唯独为你的罪过，你应常怀忧伤。不可停止劳作，即使你富足；因为懒惰之人因怠惰而多加重罪责。\newline{}\par

16. 要喜爱贫穷，渴求匮乏。若你拥有两者为业，你便是高处的继承者。不可轻视穷人的声音，也不可让他有机会诅咒你。因为若他口舌苦涩地诅咒，主必垂听他的恳求。若他的衣服污秽，用水洗净，水是无偿获得的。穷人进了你的家？天主进了你的家；天主住在你的居所。你让他从患难中得舒畅的那位，也必从患难中拯救你。你洗了客人的脚？你洗去了你罪恶的污秽。你在他面前摆设了筵席？看，天主在享用，基督也在饮用，圣神在安息。穷人因你的宴席而得饱足、得舒畅？你满足了基督你的主。他预备作你的报偿者：在天使和众人面前，他将承认你曾喂养他的饥饿；他必因你曾给他喝的、解了他的渴而感谢你。\newline{}\par

17. 哦，主是多么仁慈！哦，祂的仁慈何等无量！有福的人族，当天主宣认它时！祸哉，那个祂所否认的灵魂！火已为它的惩罚而储藏。在望德中鼓起勇气，我儿，播种美善（种子），不要灰心。农夫怀着望德播种，商人怀着望德旅行，你也喜爱美善（种子）；在望德中期待奖赏。没有祈祷的开始，什么事也不要做。以活十字架的记号，封印你所有行为，我儿。不画十字，不可出门。无论在吃或喝，无论睡或醒，无论在家或路上，或休闲之时，不可忽略这记号；因为没有任何守护者比得上它。它将成为你的一面墙，在你所有行为的前面。并将此教导给你的子女，使他们谨慎地遵从它。\par

18. 你当将自己置于法律之下，如此你才能真正成为自由人。不要随从你灵魂的欲望而背离天主的法律。我需写下多少诫命，刻下多少法律？你若渴望自由，岂非都能从自身学到？你若爱纯洁，你也会将它教给别人。让自然成为你的书，一切受造物成为你的法板；从它们学习法律，默想那未写之事。太阳在其运行中教导你休息，停止劳作。黑夜在其静默中向你呼喊，为你的工作设立了限度。大地与树上的果实呼喊，万事都有其季节。你冬天播下的种子，夏天收获其果实。同样，在世上播种正义的种子，在复活时收获它们。鸟儿每日觅食，谴责贪婪者的贪心，斥责抢夺他人积蓄的勒索。死亡，万物的终限，它自身就是万物的谴责者。\par

19. 你要投靠那不变不逝的天主。以受苦抑制欢笑，以忧愁抑制欢乐。以望德安慰受苦，以期待安慰悲伤。明智的人，你要相信并信赖，因为是你引导你的天主；若祂的眷顾不离弃你，就没有什么能伤害你。如果一个人能被另一个人拯救，卑微者能被伟人拯救，那么天主的投靠岂不更能保存那相信的人？不要因暴力攻击你的敌人而害怕。祂必警醒护卫你的灵魂，有害之物反成为有益。没有人能强制带领你，除非有自由的地方。没有人会陷入超过他力量限度的诱惑。惩罚中并无邪恶，只要自由是甘愿的。自由的行为并不乖戾，是它的意志被乖戾了。\par

20. 对于公正正直的人，考验成为帮助。约伯，一个有辨识力的人，在考验中得胜。疾病临到他，他不抱怨；病痛折磨他，他不发怨言；他的身体衰败，力气离去，但他的意志并未削弱。他因苦难而在一切事上显出完美，因为考验没有压垮他。亚巴郎从本乡、本族、本家中被召出，成为一个外方人。但这并未损害他，反而使他大大得胜。同样，若瑟从为奴之家被提升，统治埃及如王。阿纳尼亚和达尼尔的一群人从束缚中解救他人。所以，明智的人，请看自由拥有的力量：若非意志被削弱，没有什么能伤害它。以色列因奢侈生活而肥胖、踢跳，忘记了盟约。他敬拜虚无的神，忘记了受造的本性。他在沙漠的安息中忘记了埃及的束缚。每当他受磨难时，他只承认主；但当他住在安息中时，他忘记了天主，他的救赎者。不要在此寻找安息，因为这是劳苦的世界。你若能明智分辨，切莫以时间换时间；以那永存的换那不留存的；以那不止息的换那止息的；不以真理换谎言；不以身体换影子；不以警醒换沉睡；不以合时宜的换不合时宜的；不以那时间换这诸时间。要收敛你的心思，不要让它游荡于无益的多样事物中。\par

21. 在受造界中，没有敬畏天主的人是富足的；没有缺乏真理的人是真正贫穷的。那需要之人是何等贫乏，而非富足，他的需要见证了他甚至要从卑贱者和乞丐那里接受礼物。他真正是一个奴仆，且有许多主人：他服侍金钱、财富和财产。他的主人毫无仁慈，因为他们不给他安息。你要逃避，生活在贫穷中，如慈母怜悯她的爱子。投靠贫困，她以精选之物养育她的子女；她的轭是轻省甘甜的，对她的怀念是甘甜的。只有良心患病的人厌恶贫穷之饮；胆怯者畏惧那可敬的贫困之轭。人子，是谁准许你在世上找到安息？尘土之躯，是谁准许你在贫穷中富足？不要因欲望而贫乏、仰赖他人。你每日的粮食，由你脸上的汗而来，已足足够用。让这成为你需要的尺度，就是每天所给你的；若你为自己找到盛宴，从中取你所需。你不可一日取多日之粮，因为肚腹不藏珍宝。当你满足时，要赞美并感恩，免得你惹怒那施予者。要在纯洁中坚固自己，好从中获得益处。凡事要感恩赞美天主为救赎者，愿祂因祂的恩宠，使我们能听从并承行祂的旨意。\par

你们，我给了你们生命的劝告，不要在这事上疏忽。我从他人的（道理）给你们写了这些；注意不要轻视他们的话。若我在你们之前离去，要在你们的祈祷中纪念我。在每一个季节祈祷恳求，使我们的爱持续真诚。至于我们，为这些事，让我们向圣父、圣子和圣神献上赞美和荣耀，从今直到永远。阿们。\par

\SourceRecord{Translated by A. Edward Johnston.；Nicene and Post-Nicene Fathers, Second Series；Vol. 13.；Edited by Philip Schaff and Henry Wace.；Buffalo, NY: Christian Literature Publishing Co.,；1898.；<http://www.newadvent.org/fathers/3707.htm>.}
